\chapter{System Design}
\label{chap:system-design}

\section{Architecture Overview}
\label{sec:architecture-overview}

The evaluation pipeline follows a reverse-engineering design. Instead of starting from an arbitrary synthetic dataset, the system first generates a controlled real dataset with known ground truth and then synthesizes a counterpart with the CART method. This makes it possible to compare the synthetic output directly against the original geometry.

The pipeline has five stages:
\begin{enumerate}
  \item generate a real reference dataset with known cluster labels;
  \item synthesize a new dataset from the real one using CART;
  \item apply K-Means and Hierarchical Clustering to the synthetic dataset;
  \item align predicted clusters to the ground truth through the Hungarian Algorithm;
  \item compute success and structural-fidelity metrics, complemented with PCA and t-SNE visualizations.
\end{enumerate}

\section{Reference Data Generator}
\label{sec:reference-data-generator}

The real datasets are generated through controlled multivariate simulations. The design varies sample size, number of clusters, separation, dimensionality, correlation, and distribution family. This gives the study a clean baseline from which synthesis errors can be measured.

Two distributional regimes are emphasized:
\begin{itemize}
  \item \textbf{Normal datasets,} representing spherical and relatively well-behaved cluster geometry.
  \item \textbf{Gamma datasets,} representing skewed and non-spherical cluster geometry.
\end{itemize}

For the approximately Gaussian scenarios, the simulations rely on multivariate-normal generation routines available in \texttt{mvtnorm} \cite{mvtnormmanual,genzbretz2009}. For the skewed gamma scenarios, the dependence structure is imposed through copulas, implemented with the \texttt{copula} package \cite{yan2007}. This contrast is crucial because the two clustering algorithms rest on different geometric assumptions.

\section{Synthetic Data Generator}
\label{sec:synthetic-data-generator}

The synthetic component is built on the CART synthesis method \cite{nowok2016}. Given a vector of variables $X = (X_1, X_2, ..., X_p)$, the chain-rule decomposition of the joint distribution can be written as:
\begin{equation}
  f(X_1, ..., X_p) = f(X_1)\, f(X_2|X_1)\, f(X_3|X_1, X_2)\, \cdots\, f(X_p|X_1, ..., X_{p-1})
\end{equation}

In practice, the synthesis proceeds variable by variable:
\begin{enumerate}
  \item the first variable is resampled from its observed marginal behavior;
  \item a decision tree predicts the next variable from the previously generated ones;
  \item new values are sampled from the leaf-node distributions;
  \item the process repeats until all variables are synthesized.
\end{enumerate}

The factorization itself is exact; the approximation enters because each conditional distribution is fitted through CART rather than known analytically. This design is flexible and non-parametric, but it also introduces orthogonal decision boundaries. Those staircase-like artifacts are precisely what the downstream evaluation is designed to detect.

\section{Experimental Sweep}
\label{sec:experimental-sweep}

The benchmark is constructed as a Cartesian product of sample-size settings, cluster counts, separations, dimensionalities, correlation levels, distribution families, and repeated runs. In concrete terms, the sweep combines:
\begin{itemize}
  \item ten sample-size settings across the interval $N \in [100, 1000]$;
  \item three values of $k$;
  \item four separation levels;
  \item three dimensionalities;
  \item two correlation settings;
  \item two distribution families; and
  \item repeated runs for each scenario, together with both real and synthetic outputs.
\end{itemize}

This Cartesian-product design explains why the full benchmark exceeds 1,800 generated datasets.

\section{Evaluation Components}
\label{sec:evaluation-components}

The downstream system contains three logical components:
\begin{itemize}
  \item \textbf{Clustering module:} runs k-means and Ward-based Hierarchical Clustering on the synthetic datasets.
  \item \textbf{Alignment module:} solves label permutation with the Hungarian Algorithm \cite{kuhn1955} before any score is computed.
  \item \textbf{Diagnostic module:} produces numerical metrics and low-dimensional visualizations to inspect separability and structural distortion.
\end{itemize}

The combination of these components allows the report to move from a simple success-or-fail benchmark to a structural explanation of algorithmic robustness.
